\documentclass[11pt]{article}
\usepackage[margin=1in]{geometry}
\usepackage[T1]{fontenc}
\usepackage{lmodern,microtype,hyperref,enumitem,amsmath,amssymb,xcolor,booktabs}
\usepackage[numbers,sort&compress]{natbib}
\hypersetup{colorlinks=true,urlcolor=blue,citecolor=blue,linkcolor=blue}
\title{Q14: Comprehension Check}
\author{DLS 2026 --- Kukla}
\date{October 2026}
\begin{document}
\maketitle
\textbf{Comprehension check.} A student says ``I'll just find the neuron that stores the fact that Paris
is the capital of France and read it off.'' Give two reasons from this module why that is harder than it
sounds, and name where such factual knowledge is most likely concentrated.

\vspace{0.8cm}\hrule\vspace{4pt}
{\small\textbf{Answer.} (1) Concepts are represented as \emph{directions} (linear combinations of many
neurons), not single units, so there is usually no one ``Paris'' neuron to read. (2) \emph{Superposition}:
networks pack more features than neurons, so individual neurons are \emph{polysemantic} (fire for many
unrelated things), making any single neuron's activation ambiguous. Such factual associations are most
likely concentrated in the \textbf{MLP/feed-forward sublayers}, which act as key--value memories and can be
localized with causal tracing (and disentangled with sparse autoencoders).}
\end{document}
